% sec10_5.tex — Section 10.5: Application: The Demand for Money in the United States

The literature on the estimation of the money demand function is very large (see
Goldfeld and Sichel, 1990, for a review). The money demand equation typically
estimated in the literature is
\begin{equation}
  m_t - p_t = \mu + \gamma_y y_t + \gamma_R R_t + z_t^*,
  \label{eq:10.5.1} \tag{10.5.1}
\end{equation}
where $m_t$ is the log of money stock in period $t$, $p_t$ is the log of price level, $y_t$ is
log income, $R_t$ is the nominal interest rate, and $z_t^*$ is some error term. $\gamma_y$ is the
income elasticity and $\gamma_R$ is the interest semielasticity of money demand (it is a
semielasticity because the interest rate enters the money demand equation in lev\-els,
not in logs). Most empirical analyses that predate the literature on unit roots
and cointegration suffer from two drawbacks. First, as will be verified below, all
the variables involved, $m_t - p_t$, $y_t$, and $R_t$, appear to contain stochastic trends. If
the variables have trends, conventional inference under the stationarity assumption
is not valid. Second, the regressors may be endogenous. This is likely to be a seri\-ous
problem in the case of money demand, because in virtually all macro models a
shift in money demand represented by the error term $z_t^*$ affects the nominal interest
rate and perhaps income. A resolution of both problems is provided by the econo\-metric
technique for estimating cointegrating regressions presented in the previous
section.

Another issue addressed in this section is the stability of the money demand
function. The consensus in the literature is that money demand is unstable. This
is challenged by Lucas (1988) who, using annual data since 1900, argues that the
interest semielasticity is stable if the income elasticity is constrained to be unity.
We will examine Lucas' claim by applying the state-of-the-art econometric tech\-niques
developed in this chapter.

\subsection*{The Data}

We base our analysis on the annual data studied by Lucas (1988), extended by
Stock and Watson (1993) to cover 1900--1989. Their measures of the money stock,
income, and the nominal interest rate are M1, net national product (NNP), and
the six-month commercial paper rate (in percent at an annual rate), respectively.
The use of \emph{net} national product --- rather than gross national product --- follows the
tradition since Friedman (1959) that the scale variable ($y_t$) in the money demand
equation should be a measure of wealth rather than a measure of transaction vol\-ume.
There are no official statistics on M1 that date back to as early as 1900, so one
needs to do some splicing on series from multiple sources. For the money stock
and the interest rate, monthly data were averaged to obtain annual observations.
See Appendix~B of Stock and Watson (1993) for more details on data construction.

\subsection*{$(m - p, y, R)$ as a Cointegrated System}

Figures~10.2 and~10.3 plot $m - p$, $y$, and $R$. The three series have clear trends.
Log real money stock ($m - p$) grew rapidly over the first half of the century, but
experienced almost no growth thereafter until 1981. Log income ($y$), on the other
hand, grew steadily over the entire sample period with a major interruption from
1930 to the early 1940s, so M1 velocity ($y - (m - p)$), also plotted in Figure~10.3,
dropped during that period and then grew steadily until 1981. This movement in
M1 velocity is fairly closely followed by the nominal interest rate, which suggests
that $m - p$, $y$, and $R$ are cointegrated, with a cointegrating vector whose element
corresponding to $y$ is about~1.

Inspection of these figures suggests that the three variables, $m - p$, $y$, and
$R$, might be well described as being individually I(1). In fact, Stock and Watson
(1993), based on the ADF tests, report that $m - p$ and $y$ are individually I(1) with
drift, while $R$ is I(1) with no drift. They also report, based on the Stock-Watson
(1988) common trend test, that the three-variable system is cointegrated with a
cointegrating rank of~1. In what follows, we proceed under the assumption that
$m - p$ is cointegrated with $y$ and $R$. (In the empirical exercise, you will be asked
to check the validity of this premise.)

\begin{figure}[htbp]
  \centering
  % Placeholder for Figure 10.2
  \includegraphics[width=0.9\textwidth,trim=50 350 50 350,clip]{figures/fig_10_2.png}
  \caption{U.S.\ Real NNP and Real M1, in Logs}
  \label{fig:10.2}
\end{figure}

\begin{figure}[htbp]
  \centering
  % Placeholder for Figure 10.3
  \includegraphics[width=0.9\textwidth,trim=50 350 50 350,clip]{figures/fig_10_3.png}
  \caption{Log M1 Velocity (left scale) and Short-Term Commercial Paper Rate (right scale)}
  \label{fig:10.3}
\end{figure}

\subsection*{DOLS}

The first line of Table~10.2 reports the SOLS estimate of the cointegrating regres\-sion
(10.5.1). (The sample period is 1903--1987 because in the DOLS estimation
below two lead and lagged changes will be included in the augmented cointegrating
regression.) The standard errors are not reported because the asymptotic distribu\-tion
of the associated $t$-ratio, being dependent on nuisance parameters, is unknown.
The SOLS estimate of the income elasticity of 0.943 is close to unity, but we can\-not
tell whether it is insignificantly different from unity. To be able to do inference,
we turn to the DOLS estimation of the cointegrating vector, which is to estimate
the parameters $(\mu, \gamma_y, \gamma_R)$ in the cointegrating regression (10.5.1) by adding $\Delta y_t$,
$\Delta R_t$, and their leads and lags. Following Stock and Watson (1993), the number
of leads and lags is (arbitrarily) chosen to be~2, so the augmented cointegrating
regression associated with (10.5.1) is
\begin{align}
  m_t - p_t
  &= \mu + \gamma_y y_t + \gamma_R R_t
  \notag \\
  &\quad + \beta_{y0} \Delta y_t
    + \beta_{y,-1} \Delta y_{t+1} + \beta_{y,-2} \Delta y_{t+2}
    + \beta_{y1} \Delta y_{t-1} + \beta_{y2} \Delta y_{t-2}
  \notag \\
  &\quad + \beta_{R0} \Delta R_t
    + \beta_{R,-1} \Delta R_{t+1} + \beta_{R,-2} \Delta R_{t+2}
    + \beta_{R1} \Delta R_{t-1} + \beta_{R2} \Delta R_{t-2}
    + v_t.
  \label{eq:10.5.2} \tag{10.5.2}
\end{align}
This dictates the sample period to be $t = 1903, \ldots, 1987$. The DOLS point esti\-mate
of $(\gamma_y, \gamma_R)$, reported in the second line of Table~10.2, is obtained from esti\-mating
this equation by OLS.

To calculate appropriate standard errors, the long-run variance of the error term
$v_t$ needs to be estimated. For this purpose we fit an autoregressive process to the
DOLS residuals. The order of the autoregression is (again arbitrarily) set to~2. With
the DOLS residuals calculated for $t = 1903, \ldots, 1987$, the sample period for the
AR(2) estimation is for $t = 1905, \ldots, 1987$ (sample size $= 83$). The estimated
autoregression is
\begin{equation}
  \hat{v}_t = 0.93806\, \hat{v}_{t-1} - 0.13341\, \hat{v}_{t-2},
  \quad SSR = 0.19843.
  \label{eq:10.5.3} \tag{10.5.3}
\end{equation}
We then use the formula (10.4.17), but to be able to replicate the DOLS estimates
by Stock and Watson (1993, Table~III), we divide the SSR by $T - p - K$ rather
than $T - p$, where $K$ here is the number of regressors in the
augmented cointegrating regression (which is 13 here). Thus
$\hat{\sigma}_e^2 = 0.19843/(83 - 2 - 13) = 0.0029181$ and
\begin{equation}
  \hat{\lambda}_v^2
  = \frac{0.0029181}{(1 - 0.93806 + 0.13341)^2}
  = 0.07647.
  \label{eq:10.5.4} \tag{10.5.4}
\end{equation}

\begin{table}[htbp]
\centering
\caption{SOLS and DOLS Estimates, 1903--1987}
\label{tab:10.2}
\begin{tabular}{l cccc}
\toprule
 & $\gamma_y$ & $\gamma_R$ & $R^2$ & Std.\ error \\
 & & & & of regression \\
\midrule
SOLS estimates & 0.943 & $-0.082$ & 0.959 & 0.135 \\[3pt]
DOLS estimates & 0.970 & $-0.101$ & 0.982 & 0.096 \\
(Rescaled standard errors) & (0.046) & (0.013) & & \\
\bottomrule
\end{tabular}

\medskip
\small
\textsc{Source:} The point estimates and standard errors are from Table~III of Stock
and Watson (1993). The $R^2$ and standard error of regression are by author's
calculation.
\end{table}

The estimate $\hat{\lambda}_v$ is the square root of this, which is 0.277. Since the standard error
of the DOLS regression is 0.096, the factor in the formula (10.4.15) for rescal\-ing
the usual OLS standard errors from the augmented cointegrating regression is
$0.277/0.096 = 2.89$. The numbers in parentheses in Table~10.2 are the adjusted
standard errors thus calculated. Now we can see that the estimated income elastic\-ity
is not significantly different from unity.

\subsection*{Unstable Money Demand?}

Table~10.3 reports the parameter estimates by SOLS and DOLS for two subsam\-ples,
1903--1945 and 1946--1987. (Following Stock and Watson (1993), the break
date of 1946 was chosen both because of the natural break at the end of World
War~II and because it divides the full sample nearly in two.) In sharp contrast to
the estimates based on the first-half of the sample, the postwar estimates are very
different from those based on the full sample.

\begin{table}[htbp]
\centering
\caption{Money Demand in Two Subsamples}
\label{tab:10.3}
\begin{tabular}{l cccc}
\toprule
 & \multicolumn{2}{c}{1903--1945} & \multicolumn{2}{c}{1946--1987} \\
 & $\gamma_y$ & $\gamma_R$ & $\gamma_y$ & $\gamma_R$ \\
\midrule
SOLS estimates & 0.919 & $-0.085$ & 0.192 & $-0.016$ \\[3pt]
DOLS estimates & 0.887 & $-0.104$ & 0.269 & $-0.027$ \\
(Rescaled standard errors) & (0.197) & (0.038) & (0.213) & (0.025) \\
\bottomrule
\end{tabular}

\medskip
\small
\textsc{Source:} Table~III of Stock and Watson (1993).
\end{table}

Lucas (1988) argues that estimates like these are consistent with a stable
demand for money with a unitary income elasticity. To support this view, he points
to the evidence depicted in Figure~10.4. It plots $m_t - p_t - y_t$ (the log inverse velocity),
which would be the dependent variable in the money demand equation (10.5.1) if
the income elasticity $\gamma_y$ were set to unity) against the interest rate, with the postwar
observations indicated by a different set of symbols from the prewar observations.
The striking feature is that, although postwar interest rates are substantially higher,
postwar observations lie on the line defined by the prewar observations.\footnote{%
In Lucas' (1988) original plot, the sample period is 1900--1985 and the break date is 1957, but the message
of the figure is essentially the same.}

\begin{figure}[htbp]
  \centering
  % Placeholder for Figure 10.4
  \includegraphics[width=0.9\textwidth,trim=50 350 50 350,clip]{figures/fig_10_4.png}
  \caption{Plot of $m - p - y$ against $R$}
  \label{fig:10.4}
\end{figure}

This finding suggests that the postwar estimates suffer from the near multi\-collinearity
problem: because of the similar trends in $y$ and $R$ in the postwar period
(shown in Figures~10.2 and~10.3), $\gamma_y$ and $\gamma_R$ estimated on postwar data are strongly
correlated. Even though for each element the postwar estimate of $(\gamma_y, \gamma_R)$ is differ\-ent
from the prewar estimate, they may not be jointly significantly different. This
possibility could be easily checked by the Chow test of structural change, at least
if the regressors were stationary. In the Chow test, you would estimate
\begin{equation}
  m_t - p_t = \mu + \gamma_y y_t + \gamma_R R_t
    + \delta_0 D_t + \delta_y y_t D_t + \delta_R R_t D_t + z_t^*,
  \label{eq:10.5.5} \tag{10.5.5}
\end{equation}
where $D_t$ is a dummy variable whose value is 1 if $t \geq 1946$ and 0 otherwise.
Thus, the coefficients of the constant, $y_t$, and $R_t$ are $(\mu, \gamma_y, \gamma_R)$ until 1945 and
$(\mu + \delta_0, \gamma_y + \delta_y, \gamma_R + \delta_R)$ thereafter. If the regressors are stationary, the Wald
statistic for the null hypothesis that $\delta_0 = 0, \delta_y = 0, \delta_R = 0$ is asymptotically
$\chi^2(3)$ as both subsamples grow larger with the relative size held constant (see
Analytical Exercise~12 in Chapter~2). Now in the present case where regressors
are I(1) and not cointegrated, it has been shown by Hansen (1992b, Theorem~3(a))
that the Wald statistic (properly rescaled if necessary) has the same asymptotic
distribution (of $\chi^2$) if the parameters of the cointegrating regression are estimated
by an efficient method such as DOLS (which adds $\Delta y_t$, $\Delta R_t$, and their leads and
lags to the cointegrating regression (10.5.5)), the Fully Modified estimation, or the
Johansen ML.\footnote{%
We noted in the previous section that the Wald statistic is not asymptotically chi squared even after rescaling
if the hypothesis involves $\mu$. Hansen's (1992b) result, however, shows that the Wald statistic for the hypothesis
about the \emph{difference} in $\mu$ is asymptotically chi squared. Hansen (1992b) also develops tests for structural change
at unknown break dates.}
The DOLS estimates of (10.5.5) augmented with two leads and
two lags of $\Delta y$ and $\Delta R$ are shown in Table~10.4. The insignificant Wald statistic
supports Lucas' view that the money demand in the United States has been stable
in the twentieth century.

\begin{table}[htbp]
\centering
\caption{Test for Structural Change}
\label{tab:10.4}
\begin{tabular}{l ccccc c}
\toprule
 & $\gamma_y$ & $\gamma_R$ & $\delta_0$ & $\delta_y$ & $\delta_R$
 & Wald statistic \\
\midrule
DOLS estimates
  & 0.925 & $-0.090$ & 1.36 & $-0.52$ & 0.048 & 1.85 \\
(Rescaled standard errors)
  & (0.142) & (0.026) & (0.72) & (0.31) & (0.034)
  & ($p$-value $= 0.60$) \\
\bottomrule
\end{tabular}

\medskip
\small
\textsc{Source:} Author's calculation, to be verified in the empirical exercise.
\end{table}
